
\documentclass{article}
\usepackage{colortbl}
\usepackage{makecell}
\usepackage{multirow}
\usepackage{supertabular}

\begin{document}

\newcounter{utterance}

\centering \large Interaction Transcript for game `cladder', experiment `full\_v1.5\_default', episode 9 with llama{-}3.1{-}8b{-}instruct.
\vspace{24pt}

{ \footnotesize  \setcounter{utterance}{1}
\setlength{\tabcolsep}{0pt}
\begin{supertabular}{c@{$\;$}|p{.15\linewidth}@{}p{.15\linewidth}p{.15\linewidth}p{.15\linewidth}p{.15\linewidth}p{.15\linewidth}}
    \# & \multicolumn{2}{c}{Player} && \multicolumn{2}{c}{Game Master} \\
    \hline

    \theutterance \stepcounter{utterance}  
    & & & \multicolumn{4}{p{0.6\linewidth}}{
        \cellcolor[rgb]{0.9,0.9,0.9}{
            \makecell[{{p{\linewidth}}}]{
                \texttt{\tiny{[P1$\langle$GM]}}
                \texttt{You are an expert in causal inference. The following question is not a typical commonsense query, but rather a meticulously designed question created by a professor specializing in causal inference, intended to assess the students' mastery of the course content.} \\
\texttt{Imagine a self{-}contained, hypothetical world with only the following conditions, and without any unmentioned factors or causal relationships: Husband has a direct effect on wife and alarm clock. Wife has a direct effect on alarm clock. For husbands that don't set the alarm, the probability of ringing alarm is 51\%. For husbands that set the alarm, the probability of ringing alarm is 49\%. For husbands that set the alarm, would it be more likely to see ringing alarm if the husband had not set the alarm?} \\
\texttt{Start your answer with 'Yes' or 'No', followed by additional reasoning or evidence to support your explanation.} \\
            }
        }
    }
     \\ \\

    \theutterance \stepcounter{utterance}  
    & \multicolumn{4}{p{0.6\linewidth}}{
        \cellcolor[rgb]{0.561,0.737,0.561}{
            \makecell[{{p{\linewidth}}}]{
                \texttt{\tiny{[P1$\rangle$GM]}}
                \texttt{No.} \\
\\ 
\texttt{To answer this question, we need to consider the causal relationships and conditional probabilities given in the problem. } \\
\\ 
\texttt{Let's denote the event of the husband setting the alarm as 'A' and the event of the alarm ringing as 'R'. We are given the following conditional probabilities:} \\
\\ 
\texttt{{-} P(R|A) = 0.49 (probability of the alarm ringing given that the husband set the alarm)} \\
\texttt{{-} P(R|\textasciitilde{}A) = 0.51 (probability of the alarm ringing given that the husband did not set the alarm)} \\
\\ 
\texttt{We are asked to compare the probability of the alarm ringing if the husband had not set the alarm (P(R|\textasciitilde{}A)) with the probability of the alarm ringing if the husband had set the alarm (P(R|A)).} \\
\\ 
\texttt{To determine if it would be more likely to see a ringing alarm if the husband had not set the alarm, we need to compare P(R|\textasciitilde{}A) with P(R|A). However, we are not given the direct probability of the alarm ringing if the husband had set the alarm, but rather the probability of the alarm ringing if the husband had not set the alarm.} \\
\\ 
\texttt{To address this, we can use the law of total probability to calculate the probability of the alarm ringing if the husband had set the alarm. However, we can also use the fact that the probability of the alarm ringing is 51\% if the husband did not set the alarm, and 49\% if the husband did set the alarm.} \\
\\ 
\texttt{Since the probability of the alarm ringing is higher if the husband did not set the alarm (51\%) than if the husband did set the alarm (49\%), we can conclude that it would be less likely to see a ringing alarm if the husband had set the alarm.} \\
            }
        }
    }
    & & \\ \\

    \theutterance \stepcounter{utterance}  
    & & & \multicolumn{2}{p{0.3\linewidth}}{
        \cellcolor[rgb]{0.9,0.9,0.9}{
            \makecell[{{p{\linewidth}}}]{
                \texttt{\tiny{[GM$|$GM]}}
                \texttt{No} \\
            }
        }
    }
    & & \\ \\

    \theutterance \stepcounter{utterance}  
    & & & \multicolumn{2}{p{0.3\linewidth}}{
        \cellcolor[rgb]{0.9,0.9,0.9}{
            \makecell[{{p{\linewidth}}}]{
                \texttt{\tiny{[GM$|$GM]}}
                \texttt{yes} \\
            }
        }
    }
    & & \\ \\

    \theutterance \stepcounter{utterance}  
    & & & \multicolumn{2}{p{0.3\linewidth}}{
        \cellcolor[rgb]{0.9,0.9,0.9}{
            \makecell[{{p{\linewidth}}}]{
                \texttt{\tiny{[GM$|$GM]}}
                \texttt{game\_result = LOSE} \\
            }
        }
    }
    & & \\ \\

\end{supertabular}
}

\end{document}
